\documentclass[12pt]{article}
\usepackage[utf8]{inputenc}
\usepackage[T1]{fontenc}
\usepackage{hyperref}
\title{S\~ao Paulo Kaing\'ang Language Resources for K-12 Curriculum Development}
\author{Compiled by Claude Sonnet 5.5 (claude-sonnet-5-5) for Andr\'as Kornai}
\date{}
\begin{document}
\maketitle

\section{Introduction}
S\~ao Paulo Kaing\'ang (ISO 639-3 zkp; Glottocode saop1235) is an extinct Je-family language of Brazil. It was formerly spoken in S\~ao Paulo state along the Paranapanema River. Wikipedia states that it differed markedly from Kaingang proper and may have been influenced by Terena. My searches found no dedicated dictionary, primer, or school programme for this variety, so there is no native-language curriculum material I could cite. Any revitalization work would have to be reconstructive and draw on historical word lists and on the living Kaingang language.

\section{Primary Educational Resources}
I found no primers, textbooks, or school programmes specific to S\~ao Paulo Kaing\'ang.

\section{Linguistic and Reference Materials}
\begin{itemize}
\item Wikipedia, S\~ao Paulo Kaing\'ang language: a short overview of status, location, and dialect names. \url{https://en.wikipedia.org/wiki/S%C3%A3o_Paulo_Kaing%C3%A1ng_language}
\item Glottolog entry for S\~ao Paulo Kaing\'ang (saop1235), with classification and bibliography. \url{https://glottolog.org/resource/languoid/id/saop1235}
\item Wikipedia, Kaingang language: the closely related living language, which may serve as a comparative base. \url{https://en.wikipedia.org/wiki/Kaingang_language}
\item Lexvo entry for zkp. \url{https://finto.fi/lexvo/en/page/zkp}
\item A paper in \emph{Liames} (Unicamp) that compares Kaingang field data with older word lists. It covers Kaingang generally, and I could not confirm that it contains S\~ao Paulo vocabulary. \url{https://periodicos.sbu.unicamp.br/ojs/index.php/liames/article/view/1426}
\end{itemize}

\section{Religious and Community Resources}
I found none for this variety.

\section{Digital and Audio Resources}
I found no audio or digital resources for S\~ao Paulo Kaing\'ang.

\section{Curriculum Development Context}
The language is extinct and no speakers are known. A school curriculum would be heritage-awareness oriented. It could use the documentary record and the living Kaingang language as a reference, and any reconstructed material should be clearly marked as such. Further work would need archival research, for example in the sources cited by Wikipedia and Glottolog.
\end{document}